% Page 063 — page imprimée 60
% MEYRUEIS DE 1942 à 1944
% Notes prises par Marceline Causse de Pradines

\begin{center}
{\bfseries\Large MEYRUEIS DE 1942 à 1944}

\vspace{0.5em}
{\bfseries Notes prises par Marceline Causse de Pradines}
\end{center}

\medskip

\noindent{\bfseries Après le combat d’Aire de Côte le 1\textsuperscript{er} Juillet 1943}

Un rescapé Marcel Chacon de Nimes arrive à Meyrueis chez le Pasteur Frank Robert qui lui
fait une fausse carte d’identité comme jardinier du Château d’Ayres, lieu où il a travaillé sans
descendre à Meyrueis.

\medskip

\noindent\textit{Quelques mois plus tôt (?) } le pasteur Robert rappelle avoir camouflé à Meyrueis un jeune
israélite Diamant (était-ce son nom de guerre ?), fils d’un bijoutier d’Anvers, qui s’est laissé
attraper par la gendarmerie. Qu’est-il devenu ?

Il en fut de même pour un autre juif, Hamburger, horloger à Meyrueis, qui fut saisi par deux
allemands et enfourné dans une auto.

\medskip

\noindent{\bfseries Les jeunes des Chantiers de Jeunesse (groupement 19)}

Frank Robert voit à la poste 150 jeunes qui téléphonent à leurs familles qu’ils partent au STO
(Service du Travail Obligatoire) en Allemagne et les encourage à ne pas y aller. Il apprend
plus tard que ces 150 sont arrivés seulement 13 à Grenoble ; le train arrivé dans la plaine
languedocienne s’arrêtait : s’évadait qui voulait.

\medskip

\noindent{\bfseries Note de Frank Robert, pasteur de Meyrueis en 1944}

Meyrueis avait un chef du ravitaillement comme chef lieu de canton désigné par Vichy
Mr. Arbousset.. Frank Robert apprend par ce chef qu’un camion chargé de fromages
descendra du Causse et sera dirigé sur Millau pour l’Allemagne. Frank Robert monte alerter le
maquis de l’Espérou qui arrive et oblige le conducteur du camion à monter à l’Espérou, ce
qu’il fit ; et le ravitaillement fut ensuite distribué au maquis et à Meyrueis.

\medskip

\noindent{\bfseries Combat de la Parade (dimanche 28 Mai 1944, jour de Pentecôte)}

A Jontanels se trouvait Karl déserteur allemand et lieutenant, blessé, chef du groupe des
déserteurs allemands du maquis Bir Hakeim, transporté par M. Rolland de Meyrueis. Le Dr.
Bouisset l’a plâtré. Il fut ensuite caché sur les hauteurs, dans la nature, au dessus de Jontanels
et dans une grotte sur le chemin de Cabrillac et revint sur l’Espérou.

\medskip

Le maire de Gatuzières fit une carte d’identité pour Arjona du maquis Bir Hakeim afin qu’il
se rende auprès de Karl et qu’il recherche des blessés que Rolland est allé guérir dans la vallée
de la Jonte (\textit{Il s’agissait de Damiani et de Fages dit Fafa de Meyrueis.} ) Ils passèrent
plusieurs jours au moulin d’Ayres, soignés par le Dr. Juliette Cabanel et ravitaillés par Paul
Loupiac.
